\documentclass[10pt]{article}
    \usepackage[left=1in, right=1in, top=1in, bottom=1in, nohead, nofoot]{geometry}
    \usepackage[utf8]{inputenc}
    \usepackage[english]{babel}
    \usepackage[document]{ragged2e}
    \usepackage{datetime} 
    \usepackage{amssymb}
    \usepackage{graphicx}
    \usepackage{cancel}
    \usepackage{hyperref}
    \pagenumbering{gobble}
\begin{document}
    \begin{center} 
        \large\textbf{ECI 149 HW 1}

        \normalsize Charles Reid Hammond \\ \today

    \end{center} 

    \setlength{\parskip}{\baselineskip}

    \begin{flushleft}

        \textbf{Problem 1:} 
        
        A primary pollutant is one that is emitted directly to the atmosphere, a secondary pollutant is one that is formed in the atmosphere as a result of reactions between other compounds. $SO_2$ is a primary pollutant, while $O_3$ is a secondary pollutant formed from VOCs and $NO_2$.

        \hrulefill

        \textbf{Problem 2:} 
        
        \textit{(i)} Given the figures and assumptions listed in the HW assignment,

        $\frac{0.05 g-PM_{2.5}}{kg-fuel}\times\frac{14,435 mi}{car-yr}\times121,083\:cars\times\frac{1 \:gal-fuel}{24.7mi}\times\frac{3,785.4 cm^3}{gal}\times\frac{0.72 g}{cm^3}\times\frac{1kg}{1000g}=$ 9,643.1 kg PM$_{2.5}/yr$


        \textit{(ii):} Given the figures and assumptions listed in the HW assignment,

        $\frac{9,643.1 kg-PM_{2.5}}{yr}\times\frac{1000g}{kg}\times\frac{1}{1,024mi^2}\times\frac{1mi^2}{2.59\times10^6 m^2}=$ 3.6 $\times 10^{-3}$ g-PM$_{2.5}$/yr/$m^2$

\hrulefill


        \textbf{Problem 3:}


        \textit{Answer:} 
        
        $\frac{13.1 g-PM_{2.5}}{kg-fuel}\times\frac{100Mg-fuel}{hec}\times\frac{1 hec}{2.47 ac}\times153,336ac \times\frac{1000kg}{Mg}\times 50\%=$ 4.1$\times 10^7$ kg-PM$_{2.5}$

        So, with the given assumptions, around 10 times as much PM$_{2.5}$ was released during the Camp Fire than by light-duty vehicles during the entire year in California.

        \hrulefill

        \textbf{Problem 4:}

        $\frac{489,000\times35\times0.07g/mi+239,511,000\times0.07g/mi}{250,000,000}=$0.08 g/mi (taking into account the maximum of 2 significant figures).

        \hrulefill

        \textbf{Problem 5:}

        a)
        According to the EPA site you linked, 3,454,675.2 tons of $NO_X$ were emitted from ``fuel combustion'' and 7,558,684.7 tons of $NO_X$ were emitted from ``mobile'' sources nationally out of a total of 12,595,525.6 tons, which makes ``fuel combustion'' sources responsible for 27\% and ``mobile'' sources responsible for 60\% of total $NO_X$ emissions.  

        b) For California: 73,199.5 tons of $NO_X$ were emitted from ``fuel combustion'' and 446,620 tons of $NO_X$ were emitted from ``mobile'' sources nationally out of a total of 583,581.0 tons, which makes ``fuel combustion'' sources responsible for 12.5\% and ``mobile'' sources responsible for 76.5\% of total $NO_X$ emissions. 

        c) For California, if fires are excluded as a source: 73,199.5 tons of $NO_X$ were emitted from ``fuel combustion'' and 446,620 tons of $NO_X$ were emitted from ``mobile'' sources nationally out of a total of 546,495.0 tons, which makes ``fuel combustion'' sources responsible for 13.4\% and ``mobile'' sources responsible for 81.7\% of total $NO_X$ emissions.
        
        d) The fuel combustion category includes electricity generation and ``other''.
        
        \hrulefill

        \textbf{Problem 6:}

        a) Seven: CO, O$_3$, Pb, NO$_2$, SO$_2$, PM$_{2.5}$, PM$_{10}$.

        b) 
        
        CO: National, on-road mobile (37.2\%). California, wildfires (57.8\%).

        O$_3$ (using VOCs as proxy): National, Industrial and other processes (57.9\%). California, wildfires (51\%).

        Pb: National, non-road mobile (63\%). California, non-road mobile (86\%).

        NO$_2$: National, on-road mobile (38.7\%). California, on-road mobile (47\%).

        SO$_2$: National, fuel combustion [electricity] (69.4\%). California, wildfires (44.9\%).

        PM$_{2.5}$: National, industrial and other processes (63.4\%). California, wildfires (64.6\%).

        PM$_{10}$: National, industrial and other processes (86.8\%). California, industrial and other processes (48.1\%).

        \hrulefill

        \textbf{Problem 7:}

        $O_3$, according to the EPA in 
        \href{https://www.epa.gov/ground-level-ozone-pollution/ecosystem-effects-ozone-pollution}{\textit{Ground-level Ozone Pollution}}, can cause lung damage in humans and can reduce photosynthesis in sensitive plants.

        PM$_{2.5}$, according to the EPA in \href{https://www.epa.gov/pm-pollution/health-and-environmental-effects-particulate-matter-pm}{\textit{Particulate Matter (PM) Pollution}}, can aggravate asthma.

        \hrulefill

        \textbf{Problem 8:}

        To emit sulfur dioxide the fuel source must contain sulfur. Stationary electricity generating plants that use coal, which can contain high levels of sulfur, are responsible for most of the sulfur dioxide emissions. CO, on the other hand, results from incomplete combustion of a fuel, and when gasoline engines operate inefficiently they produce large quantities of CO. The main fossil fuels used for mobile transportation do not contain high quantities of sulfur (diesel is now required to be ultra-low-sulfur).

        \hrulefill

        \textbf{Problem 9:}

        (1) Commercial and residential energy sources, (2), agricultural sources, (3) traffic related sources.

        \hrulefill

        \textbf{Problem 10:}

        If gasoline is 84.2\% carbon by weight, and we assume gasoline is composed entirely of alkanes, then it follows that gasoline is 15.8\% hydrogen by weight. From this:

        $\frac{0.842 g-C}{0.158 g-H}\times \frac{1 mol-C}{12 g-C} \times \frac{1 g-H}{1 mol-H}=\frac{0.07 mol-C}{0.158 mol-H}$

        And since we know, that alkanes are molecules that correspond to the formula $C_xH_{2x+2}$,

        $\frac{0.07 mol-C}{0.158 mol-H}=\frac{x}{2x+2}$

        Solving yields x=7.77 mol-C, so the average molecule in gasoline is $\rightarrow C_{8}H_{18}$


        % \begin{figure}[h]
        %     \centering
        %       \includegraphics[width=0.5\textwidth]{FIGURE2.png}
        %       \caption{Present value cost plotted with the cost function.}
        %   \end{figure}

    \end{flushleft}

 


\end{document}